% 자동 생성 — make_si_tables.py. 손으로 고치지 말 것.
% 출처: sources/dye_panel_20261106.csv
\begin{table}[htb]
\centering
\caption{The 29-component array, in the order the three source studies number it. Five entries marked $^{a}$ are not chromogenic dyes; $^{b}$ marks a component whose stock solvent differs between the source studies.}
\label{tab:dyepanel}
\small
\begin{tabular}{rlrl}
\hline
\# & Compound & \# & Compound \\
\hline
1 & Anthracene$^{a}$ & 16 & Bromophenol Blue \\
2 & Pyrene$^{a}$ & 17 & Neutral Red \\
3 & Allura Red AC & 18 & Eosin Y \\
4 & Quinine & 19 & Indigo Carmine \\
5 & Rhodamine B & 20 & Phenol Red \\
6 & Methyl Blue & 21 & Crystal Violet \\
7 & Fluorescein & 22 & L-Glutathione (reduced)$^{a}$ \\
8 & Methyl Orange & 23 & 2,7-Dibromofluorene \\
9 & Nile Red & 24 & Ethyl viologen dibromide \\
10 & Safranin O & 25 & 4,4'-Dihydroxybenzophenone \\
11 & Toluidine Blue O & 26 & 2,4-Dinitrophenylhydrazine$^{a}$ \\
12 & Nile Blue A & 27 & 2,5-Diaminobenzene-1,4-dithiol dihydrochloride \\
13 & Cresol Red & 28 & 2-Hydrazinobenzothiazole$^{a}$ \\
14 & Eriochrome Black T & 29 & Rhodamine 6G \\
15 & Quinoline Yellow$^{b}$ &  &  \\
\hline
\multicolumn{4}{p{13cm}}{\footnotesize $^{a}$ 1 Anthracene: fluorophore rather than a chromogenic dye}\\
\multicolumn{4}{p{13cm}}{\footnotesize $^{a}$ 2 Pyrene: fluorophore rather than a chromogenic dye}\\
\multicolumn{4}{p{13cm}}{\footnotesize $^{b}$ 15 Quinoline Yellow: stock solvent differs between sources: DMSO in the first and second runs, water in the third. Reported as published; not resolved here.}\\
\multicolumn{4}{p{13cm}}{\footnotesize $^{a}$ 22 L-Glutathione (reduced): thiol reducing agent rather than a dye}\\
\multicolumn{4}{p{13cm}}{\footnotesize $^{a}$ 26 2,4-Dinitrophenylhydrazine: carbonyl derivatising reagent rather than a dye}\\
\multicolumn{4}{p{13cm}}{\footnotesize $^{a}$ 28 2-Hydrazinobenzothiazole: carbonyl derivatising reagent rather than a dye}\\
\end{tabular}
\end{table}
